%; whizzy chapter -dvi
% -initex iniptex -latex platex -format platex -bibtex jbibtex -fmt fmt
% 以上 whizzytex を使用する場合の設定。

%     Tokyo Debian Meeting resources
%     Copyright (C) 2012 Junichi Uekawa
%     Copyright (C) 2011, 2015, 2020 Nobuhiro Iwamatsu

%     This program is free software; you can redistribute it and/or modify
%     it under the terms of the GNU General Public License as published by
%     the Free Software Foundation; either version 2 of the License, or
%     (at your option) any later version.

%     This program is distributed in the hope that it will be useful,
%     but WITHOUT ANY WARRANTY; without even the implied warranty of
%     MERCHANTABILITY or FITNESS FOR A PARTICULAR PURPOSE.  See the
%     GNU General Public License for more details.

%     You should have received a copy of the GNU General Public License
%     along with this program; if not, write to the Free Software
%     Foundation, Inc., 51 Franklin St, Fifth Floor, Boston, MA  02110-1301 USA

%  preview (shell-command (concat "evince " (replace-regexp-in-string "tex$" "pdf"(buffer-file-name)) "&"))

%%ここからヘッダ開始。

\documentclass[mingoth,a4paper]{jsarticle}
\usepackage{monthlyreport}
% 日付を定義する、毎月変わります。
\newcommand{\debmtgyear}{2022}
\newcommand{\debmtgmonth}{5}
\newcommand{\debmtgdate}{21}
% started from zero:
% (let ((year 2021) (month 1)) (+ (* (- year 2005) 12) month -1))
% and add 1
\newcommand{\debmtgnumber}{209}

% Needed to import pandoc-generated LaTeX documents.
% See https://stackoverflow.com/questions/40438037/tightlist-error-using-pandoc-with-markdown
\providecommand{\tightlist}{%
  \setlength{\itemsep}{0pt}\setlength{\parskip}{0pt}}

\begin{document}

\begin{titlepage}
\thispagestyle{empty}
% タイトルページ:編集必要な部分は最初のマクロに飛ばすこと

\vspace*{-2cm}
第\debmtgnumber{}回 東京エリア Debian 勉強会資料\\
\hspace*{-2cm}
\includegraphics{image-assets/dotdeb.pdf}\\
\hfill{}\debmtgyear{}年\debmtgmonth{}月\debmtgdate{}日

% ここはアップデートすること
% 全角文字にしないとフォントのサイズが合わないので注意
\rotatebox{10}{\fontsize{30}{30} {\gt ｘｘｘ特集}}\\

\vspace*{-2cm}
\hfill{}\includegraphics[height=6cm]{image-assets/openlogo-nd.eps}
\end{titlepage}

\newpage

\begin{minipage}[b]{0.2\hsize}
 \definecolor{titleback}{gray}{0.9}
 \colorbox{titleback}{\rotatebox{90}{\fontsize{80}{80} {\gt デビアン勉強会} }}
\end{minipage}
\begin{minipage}[b]{0.8\hsize}
\hrule
\vspace{2mm}
\hrule
\begin{multicols}{2}
\tableofcontents
\end{multicols}
\vspace{2mm}
\hrule
\end{minipage}

\dancersection{最近のDebian関連のミーティング報告}{杉本 典充}

\subsection{2022 年 4 月度 東京エリア・関西合同Debian勉強会}

2022 年 4 月 16 日 (土) に東京エリア Debian 勉強会と関西 Debian 勉強会の合同で
オンラインによる Debian 勉強会を開催しました。参加者は 8 名でした。

% セミナー
セミナーは dictoss さんの「DDTP及びDDTSSの紹介」の発表を行い、
参加者全員でBoFを行い「Debianのおすすめポイント」の意見を出し合いました。

勉強会の終了後、参加者同士で Debian や OSS に関する話の情報交換を行いました。


\dancersection{事前課題}{杉本 典充}

今回の事前課題は以下です。

\begin{enumerate}
  \item xxx
\end{enumerate}

%この課題に対して提出いただいた内容は以下です。

\begin{multicols}{2}
{\small
\begin{prework}{ ����(yy\_y\_ja\_jp) }
���ޤ���������ǤϤʤ��Ǥ���... �긵�Υ�åץȥåפǤ� /boot �ˤ� ext2
 ����¾�ˤ� ext3���Ƕ�ȤäƤ���ǥ����ȥåפǤ� /boot �ˤ� ext3 ����¾
 �ˤ� LVM ��� ext4 ��ȤäƤ��ޤ���
\end{prework}

\begin{prework}{ �����ϥ� }
�ǥե���Ȥ�ext3����
�����ǥե���ȤΤޤޡ�
(NTFS���VM���᡼�����ext3�⤢�뤱��)
\end{prework}

\begin{prework}{ yos.takahashi }
ext3/4���˻ȤäƤޤ���ext3�Υǡ�������١����ˤĤ�������Linux2011ǯ1���˼�ɮ���ޤ�����
\end{prework}

\begin{prework}{ MATOHARA }
����inode �ϳ���������Ƥ���inode ��ưŪ�˳�����Ƥ���XFS �����򤹤뤳
 �Ȥ�¿���Ǥ���NILFS �Ͼ�����Ƥߤ��ΤǤ�����mount ���˰ʲ��Τ褦�ʥ��
 ���������ФƤޤ��ݤ��ʤȻפ��ޤ�����
\begin{commandline}
$ sudo mount /dev/sdb1 /mnt
mount.nilfs2: WARNING! - The NILFS on-disk format may change at any time.
mount.nilfs2: WARNING! - Do not place critical data on a NILFS filesystem. 
\end{commandline}
����¾NotePC �Ǥ�dm-crypt �ξ�˥ե����륷���ƥ���֤��ưŹ沽�����ꡢ
 eCryptfs �ǰŹ沽�����ꤷ�Ƥ��ޤ����񤭹��߻���CPU �򤫤ʤ���񤷤ޤ��ġ�
\end{prework}

\begin{prework}{ ��ޤ� }
��ext2-$>$reiserfs-$>$jfs-$>$xfs-$>$reiserfs-$>$ext3�ȻȤäƤ��ޤ�����

����:
\begin{itemize}
 \item reiserfs: ���ե������¿���ե�����Υ������������ӥ��Ӥ��Ƥ��ɤ���
       �����������դȤ����ⵤ�������θ�μ�žȬ�ݥ����ɤˡ�����
 \item jfs: ����v1.0��̾��ä�IBM���ꥨ�ʥ��������ˤ�xfs��ƨ��
 \item xfs: fsck==true�˴�ư����������ǯ�Ȥä���ΤΥޥ�����Ĵ����0byte
       �ե��������������Ѥ���줺ƨ����
\end{itemize}
������reiser4��Ķ���Ԥ��뤦���ˤ��줬�����ʤäơ����ext3�˸��경��htree�����ä������⤦Ŵ�Ĥʤ鲿�Ǥ⤤���Ǥ��������Ȥ����Ĥ�nilfs�ʤɤ˼��Ф��Ƥ��ޤ���ext3�������noatime���٤Ǥ����������aufs���碌�Ƽ�ʬ���Ȥ��Ȥ� *strap �Ķ��򥯥����˥󥰤��Ƥ��ޤ����¸����ƥ��Ȥ������Ǥ�����Ǥ���USB�����ư�Ǥ�ͭ�ѡ�

���LVM�ǤϤʤ�MD��Ȥäƾ�Ĺ���ܥХå����åפ򤷤Ƥ��ޤ���
 MD(sda,sdb,sdc)�ǹ��ۤ����̾��MD(����)�Dz�ư���Хå����åפλ���
 attach/detach�򤹤롣�֥��å���٥�ʤΤǥꥫ�Х��FSǤ���Ǥ��������̤�
 �������¾����ˡ���ʤ�������
\end{prework}

\begin{prework}{ henrich }
�ȤäƤ���֤�NTFS��Ĺ���󤸤�ʤ��Ǥ����͡����졣
�����Ρ����Ѥο������ǥ�������ext4�ǥե����ޥåȤ��ޤ����������㤤��Ƚ��ʤ��Ȥ����������Ƥ��ޤ���
\end{prework}

\begin{prework}{ emasaka }
�Ĥ뤷��FS��ȤäƤޤ�
\end{prework}

\begin{prework}{ �ܾ� }
ext3����Ѥ��Ƥ��ޤ����ä��Ѥ�ä����ȤϤ��Ƥ��ޤ���
FS����ʤ��Ǥ������Ƕ�Lenny��2TB��HDD��Ȥä���parted�äƤλȤ��ƶä��ޤ�����
\end{prework}

\begin{prework}{ ����@������ }
���������Ū�˳��Ѥ��Ƥ���ե����륷���ƥ��ReiserFS�Ǥ���
�Ż��ǻȤäƤ�Ķ���ext3�Ǥ�����ext3���Ÿ������ǥ��㡼�ʥ뤬
����Ʋ��Ǥ������ηи��ʸŤ������ͥ�Ǥ���...�ˤ����ꡢ
���ޤ꿮�Ѥ��Ƥ��ޤ���
�����ReiserFS�Ķ��ǤϤ���ޤǤνꤤ���ʤ��Ÿ����ڤä��ꡢ
�Ƥ�HDD�����줫�����ꤷ�Ƥ��ﳲ����ä��и���̵���Τǡ�
��³Ū�˻ȤäƤ��ޤ���
��ǯ����ReiserFS�Υᥤ��ȯ��(Hans Reiser)�����ᤵ��Ƥ��ޤ������ƥʥ󥹤��ۤ��Ƥ��ޤ�����
�����������θ��ReiserFS��¾�γ�ȯ�Ԥˤ���³�����ݼ餵��Ƥ���Τǡ��¿����ޤ�����
\end{prework}

\begin{prework}{ nozzy123nozzy }
\begin{enumerate}
 \item LVM�ˤĤ��Ƥϡ�CentOS5.5��Ƴ�������Τ��Τޤޤ����Ѥ��Ƥޤ�����
       ���������ƥब��äƤ���Volume̾�ϥǥե���Ȥ�����ȡʼºݤˤ�
       kickstart�ˤơˤ��ѹ����ƻȤäƤޤ����ʾ㳲���Υ���١����˺��뤿
       ���
\item ext3�ˤĤ��Ƥϡ�debian-sid�����Τޤ޻��ꤷ�Ƥ����Τ򤽤Τޤ޻Ȥ�
      �Ƥ����ꤷ�ޤ���relatime, noatime ���餤�Ͼ����ɲä��Ƥߤ����ʡ���
      �ϻפäƤޤ���
\end{enumerate}
\end{prework}

\begin{prework}{ �ޤ��������ؤ� }
\begin{itemize}
 \item Debian�Ǥ��ä˶Ťä����Ȥ�����ext3��ȤäƤޤ������ۥޥ����qcow2
       ���᡼���ǥ��������ѻ��ʳ��ϡ�LVM�ϻȤäƤޤ���
 \item �����ǰ����ü�ʤΤϡ��������DHCP�������Ѥ�Armadillo-J�ǻȤäƤ�
       ��JFFS�Ǥ����ǥե���ȤΥե����०�����Ǥϥ�֡��Ȥ�����������
       �ƽ��������Ƥ��ޤ��Τǡ�RAM�ΰ�˽񤭤��ߡ��Ÿ��ڤäƤ�ä��ʤ�
       ���������Ǥ���Debian��udhcp�Υ������ѥå���������ӥ�ɤ��ƻȤäƤޤ���
       \footnote{\url{http://d.hatena.ne.jp/mkouhei/20080601/1212330630}}
 \item ��Debian���ߤǡ���ʬ����ǰ��֥ۥåȤʤΤ�palm webOS�Ǥ��������Ubuntu��
       �������ޥ���������Τ餷���ΤǤ�����/etc/mtab�򸫤��35�Ԥ⤢�ꡢ
       ���ʤ����֤ʹ����Ǥ��͡�
\end{itemize}
\end{prework}
}
\end{multicols}

%\dancersection{Debian Trivia Quiz}{username}
%
%Debianの昨今の話題についてのQuizです。
%
%今回の出題範囲は\url{debian-devel-announce@lists.debian.org} や \url{debian-news@lists.debian.org}などに投稿された内容からです。
%
%\begin{multicols}{2}
%%; whizzy-master ../debianmeetingresume201211.tex
% $B0J>e$N@_Dj$r$7$F$$$k$?$a!"$3$N%U%!%$%k$G(B M-x whizzytex $B$9$k$H!"(Bwhizzytex$B$,MxMQ$G$-$^$9!#(B
%

\santaku
{DebConf13 $B$N3+:ECO$H3+:EF|$O!)(B}
{$BF|K\!"El5~ET(B 6$B7n(B20$BF|(B}
{$B%K%+%i%0%"(B $B%^%J%0%"(B 7$B7n(B8-14$BF|(B}
{$B%9%$%9!"%t%)!<%^%k%-%e(B 8$B7n(B11-18$BF|(B}
{3}
{$B%K%+%i%0%"$O(BDebConf12$B$N3+:ECO$G$9!#(B
DebConf13$B$O%9%$%9$N%-%c%s%WCO$G3+:E$G$9!#(B
6/20$B$O3'$5$sM=Dj$r6u$1$F$*$-$^$7$g$&!#(B}

\santaku
{$B@$3&$N(BWeb$B%5!<%P$G:G$b?M5$$N$"$k(BLinux $B%G%#%9%H%j%S%e!<%7%g%s(B(W3Techs$BD4$Y(B)$B$O!)(B}
{CentOS}
{Debian}
{Ubuntu}
{B}
{\url{http://w3techs.com/technologies/history_details/os-linux}$B$K7k2L$N%0%i%U$,$"$j$^$9!#(B
$B8=:_(B Linux $B$r;HMQ$7$F$$$k(B web $B%5!<%P$N(B 32.9\% $B$,(B Debian $B$rMxMQ$7$F$*$j!"$=$N3d9g$O8=:_$bA}2C$rB3$1$F$$$k$=$&$G$9!#(B}

\santaku
{Debian $B%+!<%M%k%A!<%`$N%a%s%P!<$G$"$j!"(Bkernel.org $B$N(B 3.2.y $B0BDjHG7ONs$N%a%s%F%J$G$b$"$k(B Ben Hutchings $B$5$s$,<!4|(B Debian $B0BDjHG$H0l=o$K=P2Y$5$l$k(B Linux $B%+!<%M%k$K(B (3.2 $B7ONs$N(B mainline $B$K$OL5$$(B) $BDI2C5!G=$,Ek:\$5$l$kM=Dj$G$"$k$H=R$Y$F$$$^$9!#(B
$BB?$/$NDI2CE@$NCf$K4^$^$l$J$$$b$N$O2?!)(B}
{PREEMPT\_RT}
{Hyper-V guest drivers$B$N6/2=(B}
{ARM64/AArch64$B%"!<%-%F%/%A%c%5%]!<%H(B}
{C}
{Hyper-V guest drivers$B$O(Bmainline kernel$B$G(B3.2$B$K$b4^$^$l$F$$$^$9$,!"$h$j2~A1$5$l$?(B3.4$B$+$i$N=$@5$,F3F~$5$l$^$9!#(B
PREEMPT\_RT$B$O%O!<%I%j%"%k%?%$%`$r<B8=$9$k$?$a$N(BPatch$B!"(B
linux-image-rt-amd64 , linux-image-rt-686-pae $B$N(Bmetapackage$B$G;HMQ$G$-$^$9!#(B
$B?7$7$$(BARM 64$B%S%C%H%"!<%-%F%/%A%c%5%]!<%H$O(Bmainline kernel 3.7$B$+$i(B}

\santaku
{Wookey$B$5$s$,%"%J%&%s%9$7$?(Balpha$BHG$N(BDebian port arm64 image$B$O!)(B}
{Debian/Ubuntu port image}
{Debian/KFreeBSD port image}
{Debian/GnuHurd port image}
{A}
{self-bootstrapp(non x86)$BBP1~$H$N$3$H$G$9!#(B\url{http://wiki.debian.org/Arm64Port}$B$G%9%F!<%?%9$,3NG'$G$-$^$9!#(B}

\santaku
{700,000$BHVL\$N%P%0$,Js9p$5$l$?F|$rEv$F$k(B700000thBugContest$B$N7k2L$,=P$^$7$?!#$=$NM=A[F|$HJs9pF|$O!)(B}
{2012/12/12$B$rM=A[$7$?(BDavidPrevot}
{$BM=A[F|(B:2013/02/04$B!"Js9pF|(B:2013/02/14}
{$BM=A[F|(B:2013/02/07$B!"Js9pF|(B:2013/02/14}
{$BM=A[F|(B:2013/02/14$B!"Js9pF|(B:2013/02/07}
{C}
{$B:G$b6a$$(B2013/02/14$B$rM=A[$7$?(BChristian Perrier$B$5$s$,Ev$F$^$7$?!#7k2L$O(B\url{http://wiki.debian.org/700000thBugContest}$B$G8x3+$5$l$F$$$^$9!#(B
$B$^$?!"(B800,000/1,000,000$BHVL\$N%P%0$,Js9p$5$l$kF|$rEv$F$k%3%s%F%9%H(B\url{http://wiki.debian.org/800000thBugContest}$B$b3+:E$5$l$F$$$^$9!#(B}

\santaku
{master.debian.org$B$,?7$7$$5!3#$K0\9T$5$l$^$7$?!#$3$l$O2?$N%5!<%P$G$7$g$&$+(B $B!)(B}
{@debian.org$B$N%a!<%k%5!<%P(B}
{$B%Q%C%1!<%8$N%^%9%?!<%5!<%P(B}
{$B%Q%C%1!<%8$N%9%]%s%5!<(B(mentor)$B$rC5$9%5!<%P(B}
{A}
{$B8E$$%5!<%P$O%G%#%9%/>c32Ey$,$"$C$?$N$G!"<wL?$HH=CG$5$l!"%G!<%?$,B;<:$9$kA0$K?7$7$$%5!<%P$K0\9T$5$l$^$7$?!#(Bftp-master.debian.org$B$O(BDebian$B$N(B official package $B%j%]%8%H%j$G$9!#%Q%C%1!<%8$N%9%]%s%5!<(B(mentor)$B$rC5$9$N$O(Bmentors.debian.net$B!#(B }

\santaku
{pbuilder$B$K(Bclang support$B$,DI2C$5$l$^$7$?!#C/$,=q$$$?%Q%C%A$G$7$g$&$+!)(B}
{Sylvestre Ledru}
{Junichi Uekawa}
{Hideki Yamane}
{C}
{Debian$B$N(BClang$B%5%]!<%H$OCe!9$H?J$s$G$$$^$9!#(B}

\santaku
{DPN - 2013$BG/(B3$B7n(B4$BF|9f$K<h$j>e$2$i$l$?F|K\$N%$%Y%s%H$O(B}
{Open Source Conference 2013 Tokyo/Spring}
{Open Source Conference 2013 Hamamatu}
{Open Source Conference 2013 Tokushima}
{A}
{\url{http://henrich-on-debian.blogspot.jp/2013/02/open-source-conference-2013-tokyospring.html} $B>\:Y$O8e$[$I!#(B}


%\end{multicols}


% % (query-replace-regexp "<.*?>" "")
% % (query-replace-regexp "^[    ]\+" "")

%-------------------------------------------------------------------------------
\dancersection{アプリにおける多言語対応の仕組み}{杉本典充}
%-------------------------------------------------------------------------------

\subsection{はじめに}

アプリケーションを国際化するには出力する文字列を各言語に翻訳して出力する必要があります。今回はプログラムの実行時に翻訳した文字列に置換する gettext とメッセージカタログについて調べてみました。

\subsection{Debian における I18N と L10N について}

Debian リファレンスに「第8章 I18N と L10N」の記事があります\footnote{\url{https://www.debian.org/doc/manuals/debian-reference/ch08.ja.html}}。この記事では I18N と L10N は以下のように説明しています。

\begin{itemize}
\item 国際化 (I18N\footnote{internationalization のこと。文字数が長いため簡略した記述として使われる。}): ソフトが複数のロケール (地域) を扱えるようにします。
\item 地域化 (L10N\footnote{localization のこと。文字数が長いため簡略した記述として使わる。}): 特定のロケール (地域) を扱えるようにします。 
\end{itemize}

また、Debian デベロッパー レファレンスに「8. 国際化と翻訳」の記事があります\footnote{\url{https://www.debian.org/doc/manuals/developers-reference/l10n.ja.html}}。この記事では「あなたが英語圏のネイティブスピーカーで他の言語を話さないとしても、国際化の問題について注意を払うことはメンテナとしてのあなたの責務です」と記述があります。そのため、Debian Developer を目指すにはプログラムの国際化についてどのような処理を行っているか知っておく必要があります。


\subsection{locale}

自分が使っているコンピュータを地域化 (L10N) した状態で動作させるには使う言語を設定する必要があります。その設定を locale (ロケール) といい、Debian や 多くの Linux システム では C ロケール (POSIX ロケールともいわれる) を用いています。Debian における locale 設定は Debian インストーラや "dpkg-reconfigure locales" コマンドで設定でき、locale を設定すると、LANG 環境変数に locale 名が設定されます。

なお、C ロケールは次の記法で地域を特定する仕様になっています。

\begin{itemize}
\item (言語コード2文字)\_(国コード2文字).(文字コード)
\end{itemize}

言語コード 2 文字は ISO-639、国コード 2 文字は ISO-3166 を用いることになっています。文字コードは Debian では国際化するために内部処理で扱う文字コードは UTF-8 とすることが一般的になっているため、"UTF-8" を指定することが多いと思います。

このルールに従い日本、アメリカ、イギリス、カナダの locale 名は以下のようになります。

\begin{table}[hbtp]
  \caption{locale の表記補法}
  \label{table:locale_display}
  \centering
  \begin{tabular}{cccc}
    \hline
    国名 & 言語コード & 国コード & locale 名 \\
    \hline
    日本 & ja & JP & ja\_JP.UTF-8 \\
    アメリカ & en & US & en\_US.UTF-8 \\
    イギリス & en & GB & en\_GB.UTF-8 \\
    カナダ & en & CA & en\_CA.UTF-8 \\
    カナダ & fr & CA & fr\_CA.UTF-8 \\
    \hline
  \end{tabular}
\end{table}

そのほか、locale 名には "C" というデフォルトのロケール設定\footnote{LANG=C と書いた方が馴染みがあるかもしれません。}も存在します。

\subsection{gettext と メッセージカタログ}

Debian では国際化及び地域化の処理を行う gettext パッケージ、po4a パッケージ、poedit パッケージなどを提供しています。

gettext パッケージはソフトウェアの中に埋め込んだ英語で記述した文字列を地域化した文字列に動的に置換して出力するライブラリを含んでいます。この gettext の仕組みを使うことで、プログラムの実行時に locale の値に従って地域化した翻訳文字列を表示することができます。


なお、地域化した文字列を保存するファイルを「メッセージカタログ」と呼びます。メッセージカタログは 3 つの形式があり、それぞれ異なる拡張子をもちます。

\begin{itemize}
\item pot ファイル
  \begin{itemize}
  \item xgettext コマンドでソースコードから翻訳が必要な文字列を抽出した po ファイルのテンプレートとなるテキストファイル
  \end{itemize}
\item po ファイル
  \begin{itemize}
  \item msginit コマンドで pot ファイルから生成する各言語の翻訳結果を保存するテキストファイル
  \end{itemize}
\item mo ファイル
  \begin{itemize}
  \item msgfmt コマンドで po ファイルをコンパイルして生成するバイナリファイルで、プログラムが参照する
  \end{itemize}
\end{itemize}


\subsection{C言語のアプリケーションを例にした翻訳処理の実装例}

\subsubsection{ディレクトリとファイルの配置}

まずは C 言語のソースコード main.c と手書きした Makefile を作成します。

\begin{commandline}
$ cd myhello
$ find .
.
./src
./src/main.c
./Makefile
\end{commandline}

例として用意した src/main.c と Makefile は以下です。

\begin{commandline}
#include <stdio.h>
#include <locale.h>
#include <libintl.h>

#define _(String) gettext(String)

#define PACKAGE_NAME "myhello"
#define LOCALE_DIR   "locale"

int main(int argc, char *argv[])
{
    char *modir = NULL;

    setlocale(LC_ALL, "");
    textdomain(PACKAGE_NAME);
    modir = bindtextdomain(PACKAGE_NAME, LOCALE_DIR);

    printf("modir: %s\n", modir);
    printf("%s\n", _("Hello world."));

    return 0;
}
\end{commandline}

\begin{commandline}
$ cat Makefile

#
# Makefile for myhello
#
PROG = myhello
CFLAGS =
LDFLAGS =

SRCS = src/main.c
OBJS = $(SRCS:.c=.o)

PO_SRCS = po/ja.po
PO_OBJS = $(PO_SRCS:.po=.mo)


.PHONY: clean pot poinitja pomergeja poinstallja

all: $(PROG)

$(PROG): $(OBJS) $(PO_OBJS)
        $(CC) $(LDFLAGS) -o $@ $(LIBS) $(OBJS)

clean:
        $(RM) $(OBJS) $(PO_OBJS) *~

pot:
        xgettext -k"_" -o $(PROG).pot $(SRCS)

poinitja:
        msginit --locale ja_JP.UTF-8 --input $myhello.pot --output-file po/ja.po

pomergeja:
         msgmerge -U po/ja.po myhello.pot

poinstallja:
        install -d locale/ja/LC_MESSAGES
        install $(PO_OBJS) locale/ja/LC_MESSAGES/${PROG}.mo

.SUFFIXES: .c .o .po .mo

.c.o:
        $(CC) $(DEBUG) $(CFLAGS) -c -o $@ $<

.po.mo:
        msgfmt -o $@ $<
\end{commandline}

\subsubsection{locale と gettext の処理}

locale 処理の関数を呼び出すために locale.h、gettext() 関数を呼び出すために libintl.h を include します。

呼び出す locale 処理の関数には以下の意味があります。

\begin{itemize}
\item setlocale(LC\_ALL, "");
  \begin{itemize}
  \item プログラムの locale を設定する
  \item 第 1 引数の LC\_ALL は、日付や金額などのロケールのカテゴリすべての値を変更する
  \item 第 2 引数は指定する locale で空文字を設定した場合はプログラム実行時の LANG 環境変数を設定する
  \end{itemize}
\item textdomain(PACKAGE\_NAME);
  \begin{itemize}
  \item プログラムのドメイン名を設定する
  \item 第 1 引数はドメイン名で、通常はプログラム名と同じ文字列を指定する (今回の場合は "myhello")
  \end{itemize}
\item modir = bindtextdomain(PACKAGE\_NAME, LOCALE\_DIR);
  \begin{itemize}
  \item 第 1 引数はドメイン名で、通常はプログラム名と同じ文字列を指定する (今回の場合は "myhello")
  \item 第 2 引数は参照するコンパイルしたメッセージカタログ (mo ファイル) を検索する起点のディレクトリを設定する\footnote{debian パッケージを作成する場合は /usr/share/locale 配下のメッセージカタログをインストール仕様になっています。}
  \item 今回の例では、ドメイン名は "myhello" のため myhello.mo を 相対パスの "locale" ディレクトリから探す、という意味になる
  \end{itemize}
\end{itemize}

翻訳した文字列に置換する処理を行う gettext の仕組みでは C 言語のマクロの \_(String) の String 部分が置換対象になります。

\begin{commandline}
    printf("%s\n", _("Hello world."));
\end{commandline}

メッセージカタログを準備していない状態でコンパイルしてプログラムを実行すると、"Hello world." と英語の文字列が出力されます。

\begin{commandline}
$ make
cc -c -o src/main.o src/main.c
cc -o myhello  src/main.o

$ ./myhello
podir: locale
Hello world.
\end{commandline}


\subsubsection{メッセージカタログの準備}

まずは xgettext コマンドを実行して pot ファイルを生成します。pot ファイルは各言語のメッセージカタログのテンプレートとなるファイルのため、Language や charset などが未定義になっています。

\begin{commandline}
$ xgettext -k"_" -o myhello.pot src/main.c
$ file myhello.pot
myhello.pot: GNU gettext message catalogue, ASCII text

$ cat myhello.pot
# SOME DESCRIPTIVE TITLE.
# Copyright (C) YEAR THE PACKAGE'S COPYRIGHT HOLDER
# This file is distributed under the same license as the PACKAGE package.
# FIRST AUTHOR <EMAIL@ADDRESS>, YEAR.
#
#, fuzzy
msgid ""
msgstr ""
"Project-Id-Version: PACKAGE VERSION\n"
"Report-Msgid-Bugs-To: \n"
"POT-Creation-Date: 2022-05-19 21:32+0900\n"
"PO-Revision-Date: YEAR-MO-DA HO:MI+ZONE\n"
"Last-Translator: FULL NAME <EMAIL@ADDRESS>\n"
"Language-Team: LANGUAGE <LL@li.org>\n"
"Language: \n"
"MIME-Version: 1.0\n"
"Content-Type: text/plain; charset=CHARSET\n"
"Content-Transfer-Encoding: 8bit\n"

#: src/main.c:19
msgid "Hello world."
msgstr ""
\end{commandline}

次に msginit コマンドを実行して各言語の po ファイルを生成します。今回は日本語用の ja.po を生成してみます。各言語の po ファイルは po ディレクトリにまとめて保存することが多いため、po ディレクトリを作成しています。

\begin{commandline}
$ mkdir po
$ msginit --locale ja_JP.UTF-8 --input myhello.pot --output-file po/ja.po
ユーザが翻訳に関するフィードバックをあなたに送ることができるように,
新しいメッセージカタログにはあなたの email アドレスを含めてください.
またこれは, 予期せぬ技術的な問題が発生した場合に管理者があなたに連絡が取れる
ようにするという目的もあります.

Is the following your email address?
  dictoss@localhost
Please confirm by pressing Return, or enter your email address.
dictoss@live.jp
https://translationproject.org/team/index.html を検索中... 完了.
Please visit your translation team's homepage at
  https://translationproject.org/team/ja.html
  https://translationproject.org/team/index.html
  https://translationproject.org/html/translators.html
  https://translationproject.org/html/welcome.html
and consider joining your translation team's mailing list
  <translation-team-ja@lists.sourceforge.net>

po/ja.po を生成.
\end{commandline}

生成した po/ja.po ファイルは以下となります。

\begin{commandline}
$ file po/ja.po
po/ja.po: GNU gettext message catalogue, UTF-8 Unicode text

$ cat po/ja.po

# Japanese translations for PACKAGE package
# PACKAGE パッケージに対する英訳.
# Copyright (C) 2022 THE PACKAGE'S COPYRIGHT HOLDER
# This file is distributed under the same license as the PACKAGE package.
# Norimitsu SUGIMOTO <dictoss@live.jp>, 2022.
#
msgid ""
msgstr ""
"Project-Id-Version: PACKAGE VERSION\n"
"Report-Msgid-Bugs-To: \n"
"POT-Creation-Date: 2022-05-20 23:37+0900\n"
"PO-Revision-Date: 2022-05-21 10:53+0900\n"
"Last-Translator: Norimitsu SUGIMOTO <dictoss@live.jp>\n"
"Language-Team: Japanese <translation-team-ja@lists.sourceforge.net>\n"
"Language: ja\n"
"MIME-Version: 1.0\n"
"Content-Type: text/plain; charset=UTF-8\n"
"Content-Transfer-Encoding: 8bit\n"
"Plural-Forms: nplurals=1; plural=0;\n"

#: src/main.c:19
msgid "Hello world."
msgstr ""
\end{commandline}

po/ja.po ファイルを編集し、原文の英語を日本語に翻訳します\footnote{vi などでファイルを編集してもよいですが、poedit パッケージに含まれる poeditor というグラフィカルな翻訳エディタを使うと便利です。}。

\begin{commandline}
#: src/main.c:19
msgid "Hello world."
msgstr "こんにちは世界。"
\end{commandline}

次に msgfmt コマンドを実行して po ファイルから mo ファイルを生成します。

\begin{commandline}
$ msgfmt -o po/ja.mo po/ja.po
$ file po/ja.mo
po/ja.mo: GNU message catalog (little endian), revision 0.0, 2 messages,
Project-Id-Version: PACKAGE VERSION '\343\201\223\343\202\223\343\201\253\343\201
\241\343\201\257\344\270\226\347\225\214\343\200\202'
\end{commandline}

この状態で myhello コマンドを実行してみても、プログラムの出力文字列は翻訳された文字列になりません。

\begin{commandline}
$ echo $LANG
LANG=ja_JP.UTF-8

$ ./myhello
podir: locale
Hello world.  
\end{commandline}

プログラムが読み込む mo ファイルのディレクトリとファイル名は bindtextdomain() の引数で決まります。今回は bindtextdomain("myhello", "locale") の引数で呼び出しているため mo ファイルはプログラム実行時のカレントディレクトリからの相対パス "locale/ja/LC\_MESSAGES/myhello.mo" を読み込みます\footnote{ja の部分は言語名になります。例えば LANG 環境変数がフランス語の場合は fr になります。}。

なお、debian でパッケージを作成してプログラムを提供する場合は mo ファイルを 以下のディレクトリ及びファイル名で配置する必要があります。

\begin{itemize}
\item /usr/share/locale/(言語コード2文字)/LC\_MESSAGES/(プログラム名).mo
\end{itemize}

\begin{commandline}
$ mkdir -p locale/ja/LC_MESSAGE
$ cp -p po/ja.mo locale/ja/LC_MESSAGES/myhello.mo  
\end{commandline}

この状態で myhello コマンドを実行してみると、プログラムの出力文字列は翻訳した日本語を出力できました。

\begin{commandline}
$ echo $LANG
LANG=ja_JP.UTF-8

$ ./myhello
podir: locale
こんにちは世界。
\end{commandline}


\subsubsection{ソースコードを変更した場合のメッセージカタログの更新}

ソースコードを変更して翻訳文字列が増える場合や減る場合、行番号がずれる場合があります。その場合は xgettext コマンドで pot ファイルを再生成し、msgmerge コマンドで既存の po ファイルと 再生成した pot ファイルをマージします。

msgmerge コマンドでマージしたときにマージに失敗した翻訳文字列は、その翻訳文字列の直前に "fuzzy" という目印がつきます。

\begin{commandline}
$ vi src/main.c
     printf("%s\n", gettext("Hello world."));
+    printf("%s\n", gettext("Hello earth."));
 
$ xgettext -k"_" -o myhello.pot src/main.c
$ msgmerge -U po/ja.po myhello.pot
... 完了.

$ diff -u po/ja.po~ po/ja.po
--- po/ja.po~   2022-05-20 23:21:29.263856024 +0900
+++ po/ja.po    2022-05-20 23:37:58.014853417 +0900
@@ -8,7 +8,7 @@
 msgstr ""
 "Project-Id-Version: myhello 0.0.1\n"
 "Report-Msgid-Bugs-To: \n"
-"POT-Creation-Date: 2022-05-20 21:57+0900\n"
+"POT-Creation-Date: 2022-05-20 23:37+0900\n"
 "PO-Revision-Date: 2022-05-19 21:43+0900\n"
 "Last-Translator: Norimitsu SUGIMOTO <dictoss@live.jp>\n"
 "Language-Team: Japanese <translation-team-ja@lists.sourceforge.net>\n"
@@ -21,3 +21,8 @@
 #: src/main.c:19
 msgid "Hello world."
 msgstr "こんにちは世界。"
+
+#: src/main.c:20
+#, fuzzy
+msgid "Hello earth."
+msgstr "こんにちは世界。"
\end{commandline}


\subsection{終わりに}

C 言語のプログラムで地域化の翻訳処理を行うにあたり、メッセージカタログの取り扱い方法を説明しました。gettext を使いこなしつつメッセージカタログの翻訳を行ってプログラムの国際化を目指しましょう。

\subsection{参考文献}

\begin{itemize}
\item Debian Wiki - Locale \url{https://wiki.debian.org/Locale}
\item Weblate - GNU gettext を使用したソフトウェアの翻訳 \\ \url{https://docs.weblate.org/ja/latest/devel/gettext.html}
\end{itemize}
      
% 冊子にするために、4の倍数にする必要がある。
% そのための調整
\dancersection{メモ}{}
\mbox{}\newpage

\vspace*{15cm}
\hrule
\vspace{2mm}
\includegraphics[width=2cm]{image-assets/openlogo-nd.eps}
\noindent \Large \bf Debian 勉強会資料\\
\noindent \normalfont \debmtgyear{}年\debmtgmonth{}月\debmtgdate{}日 \hspace{5mm}  初版第1刷発行\\
\noindent \normalfont 東京エリア Debian 勉強会 （編集・印刷・発行）\\
\hrule
\end{document}
